\section{Conclusions i Treball Futur}
\label{sec:conclusions}

\subsection{Conclusions}

El projecte Dal-i ha assolit tots els objectius plantejats inicialment. S'ha dissenyat i implementat un sistema complet que cobreix des de la interfície d'usuari fins al maquinari robòtic, passant per un pipeline de visió per computador, models d'IA generativa i una infraestructura \textit{cloud-native} escalable.

Les conclusions principals del projecte són les següents:

\begin{description}
  \item[Arquitectura de microserveis.] La separació entre Cloud Run (pipeline de visió, computacionalment intens) i Cloud Functions (operacions lleugeres de dades i parla) ha demostrat ser una decisió encertada. Permet escalar cada component de forma independent i reduir el cost quan el sistema està en repòs, ja que les Cloud Functions no generen cost si no hi ha invocacions.

  \item[Firebase Anonymous Auth.] L'autenticació anònima ha resultat ser una solució elegant per oferir un historial persistent per sessió sense forçar l'usuari a registrar-se. L'UID estable permet associar les dades a un navegador concret sense necessitar cap formulari de registre.

  \item[Traducció automàtica interna.] Processar sempre en anglès internament i traduir les entrades i sortides automàticament ha simplificat enormement el codi de processament d'estils, que no necessita conèixer cap idioma específic. El resultat és un sistema multilingüe robust amb un sobrecost de latència mínim.

  \item[Optimització TSP per a robòtica.] L'aplicació de tècniques d'optimització de rutes (Greedy Nearest Neighbour + 2-opt) al problema d'ordenació de contorns per al robot ha tingut un impacte positiu mesurable en el temps de dibuix. Aquesta connexió entre algorísmica clàssica i robòtica ha estat un dels aspectes més enriquidors del projecte.

  \item[Integració de maquinari.] La coexistència de clients diversos (navegador web amb Firebase Auth i Raspberry Pi sense autenticació) s'ha resolt de manera neta amb la dependència d'autenticació opcional, sense duplicar cap lògica de negoci.
\end{description}

\subsection{Aprenentatges}

El projecte ha permès adquirir experiència pràctica en:

\begin{itemize}
  \item Desplegament de contenidors Docker multi-etapa a Google Cloud Run.
  \item Implementació i desplegament de Cloud Functions de 2a generació amb Python.
  \item Configuració d'autenticació Firebase en aplicacions web estàtiques (Next.js static export).
  \item Integració de múltiples APIs de GCP (Speech-to-Text, TTS, Translation, Vertex AI) en un sol sistema coherent.
  \item Algorísmica aplicada (TSP) en un context de robòtica real.
  \item Gestió de problemes de CORS, IAM i dominis autoritzats en entorns de producció.
\end{itemize}

\subsection{Treball Futur}

Com a línies de treball futur, proposem les extensions següents:

\begin{enumerate}
  \item \textbf{Generació de vídeo del procés de dibuix.} Capturar la trajectòria del robot en temps real i generar un vídeo time-lapse del dibuix, que es podria compartir directament des de la interfície.

  \item \textbf{Col·laboració en temps real.} Permetre que múltiples usuaris contribueixin a un mateix dibuix simultàniament, amb sincronització via WebSockets o Firebase Realtime Database.

  \item \textbf{Ampliació d'estils artístics.} Incorporar un mecanisme de \textit{fine-tuning} de Gemini per afegir nous estils artístics de manera dinàmica, sense canvis al codi.

  \item \textbf{Dashboard d'administració.} Crear un panell de control per monitoritzar l'ús del sistema, els costos de GCP i les estadístiques de dibuixos generats.

  \item \textbf{Optimització del pipeline amb GPU.} Explorar l'execució del pipeline de visió en una instància amb GPU a Cloud Run per accelerar el processament de Mode Avançat.

  \item \textbf{Exportació en formats estàndard.} Permetre descarregar les trajectòries generades en formats industrials com G-code o SVG optimitzat, compatibles amb una gamma més àmplia de màquines CNC i plotters.
\end{enumerate}
